\setcounter{section}{13}
\section{Modelo Economico}

Nombre completo: El modelo económico español en el marco de la economía mundial. El crecimiento sostenible y las políticas ambientales. La Agenda 2030 y los Objetivos de Desarrollo Sostenible. Líneas generales de la política económica actual. La renta nacional. Evolución y distribución. La internacionalización de la economía española. Comercio exterior y balanza de pagos.

\noindent\fbox{\parbox{\textwidth}{
    Los apuntes originales se hicieron en 2020
}}

\localtableofcontents

\subsection{Principios rectores de la política económica, social y medioambiental} %5
    La Constitución detalla derechos económicos, sociales y medioambientales. Tras la Agenda 2030 se busca un equilibrio entre las tres.

    En la parte económica, los principios rectores de la política económica cumple la función de orientar a los poderes públicos, como unas líneas base para la corrección de situaciones de desequilibrio del funcionamiento ordinario.

    La constitución española reconoce el derecho a la propiedad privada. Esta puede quedar limitada por la función social. También se reconoce la libertad de empresa en la economía de mercado.

    Se obliga a los poderes públicos a promover las condiciones favorables para el progreso social y económico y para una distribución regional y personal equitativa, también se realizará una política orientada al pleno empleo.

    Los Art. 40 a 42  regulan Las actividades laborales, Regulando el derecho a la formación, seguridad, higiene y garantíá del descanso con límites de duración de jornada y vacaciones. También se regula en el Art. 41 la seguridad social, garantizando la asistencia y prestaciones ante situaciones de necesidad, con énfasis en el desempleo. Por último dicta al estado a velar por los derechos de los trabajadores en el extranjero.

    Socialmente hablando, los derechos individuales se encuentran positivizados y gozan del máximo nivel de protección, mientras que los derechos sociales aparecen como principios programáticos.

    Se establece desde el Art. 1 que Espaæa es un Estado social y democrático de derecho. Desde el Art.9 se promoverán las condiciones para la realidad efectiva de las libertades e igualdades de los individuos y los grupos que lo conforman. También máxima protección en los Art. 27 y 28 de libertad sindical y huelga. Laboralmente al trabajo y negociación colectiva en los ARt. 35 y 37.

\subsection{Renta nacional y Balanza de Pagos} %8
    El \textbf{sistema de cuentas nacionales} constituye un sistema de cuentas que reflejan las magnitudes macroeconómicas con los aspectos relevantes de la economía de un país. Siendo un resumen de las actividades bajo un mismo marco analítico. Usa nomenclatura del FMI internacionalmente aceptada. Se divide en cuatro bloques, en el ámbito doméstico están los sectores real, monetario/financiero y público/gubernamental, en el bloque exterior, el sector externo.

    La Renta Nacional, el Producto Nacional Bruto y el Producto Interior Bruto son magnitudes económicas de síntesis de los resultados anuales del sistema económico del país. Son mas importantes cuando se analiza su evolución, estando vinculada a los efectos del desarrollo, nivel de vida, y bienestar social. También inciden en el crecimiento, la prosperidad y el desarrollo.

    Generalmente los tres valores están relacionados y evolucionan de forma conjunta. Se suele hablar de productos finales para no realizar un doble cómputo de productos intermedios
    
    El \textbf{Producto Interior Bruto} se refiere a los bienes y servicios finales producidos por los residentes en territorio nacional. Se le resta el consumo de bienes inservicios intermedios usados en la producción. Está relacionado con la renta total generada en un país e incluye a los ciudadanos extranjeros residentes.

    El \textbf{Producto Nacional Bruto} son los bienes y servicios finales producidos por los nacionales de un país, independientemente de si se encuentran en el país o en el extranjero. Se suele calcular partiendo del PIB, sumando las rentas obtenidas en el extranjero y restando las de los no residentes.

    El \textbf{Producto Nacional Neto} se obtiene del bruto, restando el consumo de capital fijo (amortización).

    La \textbf{Renta Nacional} es la suma de todos los ingresos de los factores productivos nacionales

    También se pueden aplicar los valores \textit{per capita} dividiendo por la población total.

    \subsubsection{Evolución de la renta nacional}
        La Renta Nacional ha ido en consonancia con las diferentes etapas de la economía española. Cabe destacar los siguientes momentos:

        \begin{description}
            \item[1977-1986] Dificultades económicas y reformas. Crecimiento reducido y aumento del desempleo con el aumento de la población activa.
            \item[1986-1992] Entrada en la Comunidad Europea y crecimiento económico. Se facilita la apertura económica española con un incremento de la inversión extranjera. Años de gran crecimiento exonómico y reducción del paro(del 20\% al 6\%). También se modera la inflacción.
            \item[1993-1994] Gran crisis económica con alto paro. Alta inestabilidad exterior y excesos de inversión en las olimpiadas y la expo de Sevilla. El gobierno realiza 3 devaluaciones de la peseta en un año
            \item[1995-2007] Crecimiento económico y busbuja inmobiliaria. Entrada en la UEM, cumpliéndose los criterios de convergencia, lo que provocó una bajada de los tipos de interés y confianza de inversores.
            \item[2007-2014] Crisis económica, se llegó a tener crecimiento 0\% en 2010 y decrecimiento del 1\% en 2011, un 2,9\% en 2012 y un 1,7\% en 2013.
            \item[2014-2019] Recuperación de la economía y creación de empleo. Se suman años de avances, basado en las fortalezas de las exportaciones, el consumo y la inversión empresarial.
            \item[2020] Crisis COVID 
        \end{description}

    \subsubsection{Medición de la desigualdad}
        Hay múltiples medidores de desigualdad, siendo los mas importantes la curva de Lorenz, el coeficiente de Gini, la convergencia sigma y los indicadores de pobreza.

        La \textbf{curva de Lorenz} mide la diferencia entre el acumulado de la renta real e ideal. En una sociedad equitativamente distribuida esta curva acumulada es lineal, mientras que en la realidad tiene forma de stick de hockey, midiendo el área debajo de ambas curvas.

        El \textbf{coeficiente de Gini} pone números a esa diferencia. Siendo la proporción entre el área entre las dos curvas, y el área de la curva ideal. No tener desigualdad supone un coeficiente de 1.

        El \textbf{Índice de convergencia sigma} Estudia una evolución temporal de la desigualdad. En un año mide la desviación estándar de las rentas familiares. Idealmente tomaría el valor 0.

        La \textbf{Tasa de pobreza} se define como el porcentaje de familias con niveles de renta por debajo de un umbral. La elección de este umbral es relativa y subjetiva.Se suele basar en los ingresos por consumo del hogar netos (divididos por el número de unidades de consumo). Normalmente este umbral se fija en el 60\% de la mediana de ingresos por unidad de consumo.
        
    \subsubsection{Políticas de redistribución de renta}














\subsection{El modelo económico español en el marco de la economía mundial} %18
\subsection{Internacionalización de la economía española} %30
\subsection{La economía española en el contexto mundial} %31
\subsection{La agenda 2030: El desarrollo sostenible incluyendo las políticas mediambientales} %34
%43